\section{Exercise 1b: Envelope Detection}
\label{sec:lab2_ex1b}

\subsection{Objective}

Implement an envelope detector using rectification and low-pass filtering, then
compare its behaviour with coherent demodulation.

\subsection{Module Design}

The implemented signal flow is:
\begin{enumerate}
\item apply a pre-gain of $\pi$ to compensate the dominant rectification loss,
\item convert the waveform into array form for sample-wise processing,
\item perform half-wave rectification by setting negative samples to zero,
\item rebuild the waveform while preserving its timing parameters,
\item apply a Butterworth low-pass filter,
\item remove the DC component and compute a PSD of the recovered signal.
\end{enumerate}

\begin{figure}[H]
\centering
\labincludegraphics[width=\linewidth]{figures/Envelope_Devodulation_VI.png}
\caption{Block diagram of the envelope detection VI.}
\label{fig:lab2_envelope_vi}
\end{figure}

\subsection{Link to Theory}

Envelope detection works by extracting the slow variation of the AM envelope.
Its validity depends on the no-inversion condition
\begin{equation}
\mu \le 1.
\end{equation}
For $\mu > 1$, the envelope inverts and strong distortion appears after
rectification and filtering.

The use of a $\pi$ pre-gain is consistent with the approximate rectification-loss
model
\begin{equation}
\hat{e}(t)\approx \frac{1}{\pi}\left[A_c + A_m m(t)\right].
\end{equation}

\subsection{Implementation Notes}

\begin{itemize}
\item timing information must be preserved when converting between waveform and
array representations, otherwise the recovered time and frequency axes become
incorrect,
\item the low-pass filter should pass the message bandwidth while rejecting
carrier-related ripple,
\item DC removal is required because the recovered envelope contains a large mean
term related to the carrier amplitude.
\end{itemize}

\subsection{Expected Output}

For $\mu \le 1$, the detector output should track the message envelope and show a
dominant spectral component near $f_m$. As $\mu$ exceeds 1, visible distortion
should remain even after filtering.
